\chapter{Testing an identification without assuming one}
\label{ch:identification}

Many theories can share an approximate central charge and a few low
dimensions. A convincing identification must also explain multiplicities,
charges, three-point tensors, locality, and the actual noninvertible
junction data. This chapter develops several discriminating tests that are
stronger than matching decimals. Some are exact theorems about candidate
theories; their application to the chain remains conditional on a controlled
continuum dictionary. Others are explicit ultraviolet constructions that
organize the observed charges but have not yielded the required infrared
theory. Keeping these roles separate allows a failed candidate to teach us
something without being mistaken for a general no-go theorem.

\section{What must match}

A proposed identification has at least four layers. First, specify the
microscopic model: the physical $\mathcal H_3$ associator, rho-channel Hamiltonian,
regular periodic space, couplings, and twists. Second, specify the limiting
kinematics: vacuum, velocity, central charge, complete low shells, and their
charge and momentum assignments. Third, specify local dynamics: normalized
OPE coefficients and a consistent crossing/sewing structure. Finally,
identify the categorical lines with their fusion, associators, and action
on the local fields and twisted sectors.

A chiral commutant with a certain central charge addresses only part of
this list. A modular-invariant torus function addresses another part.
Neither alone specifies a local full CFT. Conversely, a lattice model with
the correct ordinary $C_3\times C_3$ charges need not possess a Haagerup
line. The irrational dimension $d=(3+\sqrt{13})/2$ and its junctions are
substantial extra data.

The central charge near two motivates comparison with two compact bosons,
orbifolds, cosets, and interacting deformations. It does not justify
assuming any one of those descriptions. The exact $c=8$ research
construction and the separate irrational chiral axis discussed in
Chapter~\ref{ch:cosets} should likewise not be relabeled as the infrared limit of
the physical chain.

\section{The simplest product-theory test}

Suppose a proposal assigns
$\phi=\sigma_A\otimes1$,
$\chi=1\otimes\sigma_B$, and
$\psi=\sigma_A\otimes\sigma_B$ in a product CFT with normalized
two-point functions. Factorization immediately gives
\[
 \Delta_\psi=\Delta_\phi+\Delta_\chi,\qquad
 C(\phi,\chi,\psi^\dagger)=1.
\]
The self cubics factorize as well, so $I_\lambda=1$.
These statements are insensitive to field phases once expressed in
absolute coefficients and the invariant $I_\lambda$.

The physical chain's finite mixed coefficient near $1.42$, nonreal
$I_\lambda$, and nonadditive gaps disagree with that elementary
dictionary at currently measured sizes. Turning this into a continuum
exclusion requires limits of those quantities. It does already explain
why recognizing two ordinary $C_3$ charges is insufficient to call the
theory a product of two familiar charged CFTs.

\section{Compact bosons and a sharp light-shell bound}
\index{Narain CFT}

A rank-two Narain theory is a unitary CFT of two compact bosons with a
momentum-winding lattice of signature $(2,2)$. A vertex operator is
labeled by $(p_L,p_R)$, with
\[
 h=|p_L|^2/2,
 \qquad\bar h=|p_R|^2/2,
\]
before adding nonnegative integer oscillator levels. The lattice pairing
\[
 B(p,q)=p_L\cdot q_L-p_R\cdot q_R
\]
is integral. Normalized vertex three-point coefficients have modulus one
when the charges sum to zero, and vanish otherwise. Cocycle phases affect
their phases, not their magnitudes.
This is the two-boson version of the circle construction in
Chapter~\ref{ch:cft}; the integral pairing encodes mutual locality
between momentum-winding vertices.

A scalar vertex has $|p_L|^2=|p_R|^2=\Delta$. For scalar dimensions below
$1/2$, oscillators are absent. Any two charges in light scalar shells obey
\[
 |B(p,q)|\leq2\sqrt{\Delta_p\Delta_q}<1.
\]
Integrality therefore forces $B(p,q)=0$. All charges in the shells span a
totally isotropic real subspace of dimension at most two. Projection to
the left momenta is injective on this subspace, because a null vector
with $p_L=0$ also has $p_R=0$. Charge addition can consequently be studied
as addition of vectors in one Euclidean plane, with each dimension shell
lying on a circle.

\begin{theorem}[Unequal light scalar shells]
Let $A,B,C$ be normalized scalar fields in a compact rank-two Narain
theory, with $0<\Delta_A,\Delta_B,\Delta_C<1/2$. If their dimensions
are not all equal, then
\[
 |C(A,B,C)|\leq1.
\]
Arbitrary degeneracy and complex mixtures within each shell are allowed.
\end{theorem}
\begin{proof}
Choose the two unequal radii as the $B,C$ inputs. Write their vertex
coefficients as $b_\mu,c_\nu$, with squared sums one. For an output
charge $-\lambda$ on the $A$ circle, set
\[
 T_\lambda=\sum_{\mu+\nu=-\lambda}
           \epsilon(\mu,\nu)b_\mu c_\nu,
 \qquad |\epsilon|=1.
\]
Two circles in the plane have at most two intersection points, so
each nonzero output has at most two input decompositions. If there
are two, write
$\mu+\nu=\mu'+\nu'=-\lambda$ and
$\delta=\mu-\mu'=\nu'-\nu$. Equal radii within each input shell
give
\[
 2\mu\cdot\delta=2\nu'\cdot\delta=|\delta|^2,
 \qquad(\mu-\nu')\cdot\delta=0.
\]
The crossed input pairs have larger output radius:
\[
 |\mu+\nu'|^2=|\lambda|^2+|\delta|^2.
\]
They are absent from the diagonal sum on the $A$ circle. The
interference term is bounded by
\[
 2|b_\mu c_\nu b_{\mu'}c_{\nu'}|
 \leq |b_\mu c_{\nu'}|^2+|b_{\mu'}c_\nu|^2.
\]
Each crossed pair is used at most once. To see this, fix $(\mu,\nu')$.
Their unequal radii imply $\mu\ne\nu'$, so $\delta$ lies on a fixed
one-dimensional perpendicular. Writing $\delta=te$ gives
$t^2=2t\mu\cdot e$, with at most one nonzero solution. Thus the
original two decompositions are uniquely recovered.

Every interference term is therefore charged to distinct input pairs
outside the allowed diagonal set. Summing proves
$\sum_\lambda|T_\lambda|^2\leq
\sum_{\mu,\nu}|b_\mu|^2|c_\nu|^2=1$.
Finally, contracting $T$ with the normalized coefficient vector of
$A$ and using Cauchy--Schwarz gives the result.\end{proof}

This proof is useful beyond its numerical bound. It treats arbitrary
shell degeneracy rather than a favored vertex basis, retains cocycle
phases through absolute inequalities, and uses integrality before
geometry. It is therefore not restricted to pure momentum states or
a special torus radius.

For three equal light dimensions the sharp bound is instead
$2/\sqrt3$. The unequal-radius uniqueness argument above no longer
applies. A short separate argument gives the correct constant. Three
nonzero equal-radius vectors adding to zero form an equilateral
triad, and a charge belongs to at most one unordered such triad:
the other two must be its rotations by $\pm120$ degrees. Thus the
allowed triples partition into disjoint three-element sets.

On one triad, after taking absolute values, the output for normalized
nonnegative vectors $b,c\in\R^3$ is
$T_i=b_jc_k+b_kc_j$, with $\{i,j,k\}=\{1,2,3\}$. Direct expansion
gives
\[
 \sum_iT_i^2
 =1+\left(\sum_ib_ic_i\right)^2-2\sum_ib_i^2c_i^2
 \leq1+\frac13\left(\sum_ib_ic_i\right)^2\leq\frac43.
\]
The first inequality uses Cauchy--Schwarz on the three products
$b_ic_i$. A final contraction with a unit vector has magnitude at
most $2/\sqrt3$. If several disjoint triads occur, their contributions
are bounded by $(2/\sqrt3)\sum_Ta_Tb_Tc_T$, where the restricted
coefficient norms have squared sums at most one. H\"older's
inequality bounds that last sum by one. Cocycle phases cannot
increase this absolute bound.

For a Hermitian field, opposite triads have equal coefficient norms.
Pairing them gives the stronger self-cubic bound $\sqrt{2/3}$:
if the norm in each member of a pair is $a_T$, then
$2\sum_Ta_T^2\leq1$ and the cubic is bounded by
$(4/\sqrt3)\sum_Ta_T^3\leq\sqrt{2/3}$. The following example
attains both constants when the corresponding vertex charges exist.

\begin{example}[An equilateral vertex triad]
Choose three equal-length scalar charges $p_1,p_2,p_3$ with
$p_1+p_2+p_3=0$, and normalize
\[
 \phi=(V_{p_1}+V_{p_2}+V_{p_3})/\sqrt3.
\]
When the compatible vertex phases are chosen coherently, its cubic
gets the six permutations of the charge-conserving triple, giving
\[
 C(\phi,\phi,\phi)=\frac6{3\sqrt3}=\frac2{\sqrt3}.
\]
The Hermitian quadrature $X=(\phi+\phi^\dagger)/\sqrt2$ has cubic
$\sqrt{2/3}$. A single normalized vertex has cubic magnitude at
most one, but a degenerate-shell mixture can have a larger cubic.
\end{example}

The chain's finite $|\lambda|$ near $1.42$ is thus a strong comparator
if its limiting fields have distinct dimensions below $1/2$. The
theorem applies to a compact rank-two parent and to untwisted fields
that lift isometrically to it. It is not by itself a theorem about
twisted fields of every orbifold or about every $c=2$ CFT.

\section{What an orbifold loophole requires}

Gauging a finite symmetry introduces twisted sectors. Their local fields
need not be ordinary sums of parent vertex operators, so one cannot
apply a parent shell bound to them without an additional argument.
Complete shell multiplicity provides one such argument in specified
classes.

Suppose the entire real scalar space at a light dimension is one
charged doublet $\phi,\phi^\dagger$, with nonzero cubic and
$\langle\phi\phi^\dagger\rangle=1$. Charge conservation gives
$\langle\phi\phi\rangle=0$. Its normalized real directions are
\[
 X_\theta=\frac{e^{i\theta}\phi+e^{-i\theta}\phi^\dagger}{\sqrt2},
 \qquad C(X_\theta,X_\theta,X_\theta)
 =\frac{\Re(e^{3i\theta}c_\phi)}{\sqrt2}.
\]
The factor $1/\sqrt2$ follows by expanding the cube: only the
$\phi^3$ and $(\phi^\dagger)^3$ terms survive. On this real
two-plane the cubic is proportional to $\Re(z^3)$. Its orthogonal
stabilizer is the six-element dihedral group $D_3$. Thus every
symmetry preserving that isolated shell acts through $D_3$.

In an ordinary $\Z_2$ orbifold, the quantum involution must act on
this plane either trivially or as a reflection. In either case there
is an invariant real direction. The invariant direction lifts to
an untwisted parent field, where the Hermitian Narain bound applies.
At a suitable cubic extremum, its cubic magnitude is
$|c_\phi|/\sqrt2$, giving
\[
 |c_\phi|\leq2/\sqrt3
\]
for dimensions below $1/2$. This allows the original charged field
to mix twisted and untwisted directions; the isolated full plane is
the essential hypothesis.

For finite linear geometric point groups, the argument extends
below dimension $2/9$: fields fixed by every quantum character
lie in commutator sectors, whose nonidentity rotations have twist
cost at least $2/9$. The following result handles a larger class
of asymmetric actions, with a slightly smaller dimension interval.

\subsection{Finite asymmetric actions and the shift subgroup}
\index{orbifold!isolated doublet bound}

\begin{proposition}[The current-preserving orbifold doublet bound]
\label{prop:asymmetric-doublet}
Gauge an effective finite group $G$ of internal symmetries of a
compact rank-two Narain theory. Assume the gauging is an ordinary
consistent bosonic unitary orbifold and that $G$ preserves a chosen
two-dimensional Heisenberg current space on each chirality, setwise.
Asymmetric lattice actions, momentum/winding phases, nonabelian
groups, and consistent ordinary discrete torsion are allowed.

If the entire real scalar eigenspace at
$0<\Delta<5/36$ is one nontrivial order-three doublet, then its
normalized charged cubic obeys
\[
 |c_\phi|\leq2/\sqrt3.
\]
\end{proposition}

The premise preserves each two-dimensional current space, not
each current separately. At an enhanced point, a symmetry can mix
a chosen Cartan current with root vertex currents; such an action
is outside this proposition unless an invariant pair of Cartans
can be chosen. We first prove the stated result before addressing
those enhanced cases.

\begin{proof}
Write the even unimodular charge lattice of signature $(2,2)$ as
$\Gamma$. A current-preserving symmetry acts by
\[
 V_p\longmapsto u_g(p)V_{M_gp},\qquad
 (J_L,J_R)\longmapsto(R_L(g)J_L,R_R(g)J_R),
\]
where $R_L,R_R\in O(2)$ and $M_g$ preserves the lattice and
its chosen positive and negative planes. We will use three
pieces of this action: its current-trivial kernel, the signs of
its two determinants, and its finite rotations.

\emph{First gauge the current-trivial subgroup.}
Let $K=\ker(R_L\times R_R)$. The currents determine both
charges, so $M_k=1$ for $k\in K$. The vertex-operator OPE
then gives $u_k(p+q)=u_k(p)u_k(q)$. Thus effectiveness makes
$K$ a finite abelian subgroup of $\Hom(\Gamma,U(1))$; it is
normal as the kernel of a group action. Lattice self-duality
identifies its characters with distinct rational shifts
$s_k\bmod\Gamma$.

There are no twisted oscillators in a $k$ sector: its charges
are $\Gamma+s_k$ before projection. Define the invariant
sublattice
\[
 \Gamma_0=\{p\in\Gamma:u_k(p)=1\text{ for every }k\in K\}.
\]
The pairing between $\Gamma$ and $K$ makes the map
$\Gamma\to\widehat K$ onto. Indeed a proper image would
have a nontrivial annihilating character, hence a nonidentity
element of $K$ acting trivially on every charge, contrary to
effectiveness. Therefore $[\Gamma:\Gamma_0]=|K|$.

Projection in each $k$ sector imposes a character condition on
the original charge $p$. Discrete torsion changes which character
occurs, but surjectivity still makes the solutions one nonempty
coset of $\Gamma_0$. The $|K|$ sectors remain disjoint because
their $s_k\bmod\Gamma$ are distinct. Their union $\Gamma'$
is additive and even integral by closure and locality of the
gauged vertex OPE. Its covolume is
\[
 \operatorname{covol}(\Gamma')
 =\frac{\operatorname{covol}(\Gamma_0)}{|K|}=1.
\]
Hence $\Gamma'$ is again even unimodular of signature $(2,2)$,
with one oscillator Fock space per charge. The $K$ orbifold is
another rank-two Narain theory. This covolume argument does not
require $K$ to be a geometric target-space translation group.

Fields of the full $G$ orbifold supported on $K$ monodromies
lift isometrically into this intermediate torus theory, subject
to the remaining $G$ equivariance. Their normalized sphere
correlators agree. On a punctured sphere the local monodromies
generate the holonomies; when these lie in $K$, the torsion
restricts to $K$, and the gauge-volume factor cancels against
vacuum normalization. No extra $\sqrt{|G|}$ multiplies a
normalized cubic. We need this subspace and its correlators,
without assuming that the remaining quotient acts geometrically.

\emph{Next bound the cost of a proper rotation.}
The two maps $\det R_L,\det R_R:G\to\{\pm1\}$ are group
characters. If both signs are positive, the charge action belongs
to the identity component of real $SO(2,2)$. Its integral group
is represented by
\[
 [SL(2,\Z)\times SL(2,\Z)]/\{(I,I),(-I,-I)\}.
\]
One can see the representation on integral two-by-two matrices
with quadratic form minus the determinant: $(A,B)$ acts by
$X\mapsto AXB^{-1}$. A finite element has finite lifts $A,B$,
since a suitable power is the same central sign in both factors.
A finite-order integral $SL(2)$ matrix has trace
$-2,-1,0,1,$ or $2$. At trace $\pm2$, finite order excludes
a nonzero Jordan part, giving $\pm I$; the other traces give
orders $3,4,6$. Its rotation fraction is thus among
\[
 0,\quad\frac12,\quad\pm\frac13,\quad
 \pm\frac14,\quad\pm\frac16.
\]
For fractions $a,b$ in the two factors, the physical chiral
fractions are $a+b,a-b$, up to orientation and interchange.
For instance, split the real two-by-two matrices into the
planes commuting and anticommuting with the standard complex
structure. Left and right rotations subtract their angles on
one plane and add them on the other.
Reducing each to $[0,1/2]$ gives the complete nonidentity list:
\begin{center}
\begin{tabular}{c|c}
Chiral rotation fractions & Scalar lower bound\\\hline
$(\nu,\nu)$, $\nu=1/2,1/3,1/4,1/6$
 & $\nu(1-\nu)\geq5/36$\\
$(1/2,0)$ & $1/4$\\
$(1/3,0)$ & $2/9$\\
$(1/2,1/6)$ & $1/4$\\
$(1/12,5/12)$ & $35/144$
\end{tabular}
\end{center}
The two real chiral bosons rotated through fraction $\theta$
have oscillator ground weight $\theta(1-\theta)/2$. This
normal-ordering shift can be checked with an exponential regulator,
without assigning values to divergent sums. Put
\[
 S_a(\varepsilon)=\sum_{n\geq0}(n+a)e^{-\varepsilon(n+a)}
 =-\partial_\varepsilon
       \frac{e^{-a\varepsilon}}{1-e^{-\varepsilon}}.
\]
Its expansion is
$S_a=\varepsilon^{-2}-(a^2-a+1/6)/2+O(\varepsilon)$.
The two twisted frequency lists start at $\theta$ and
$1-\theta$, so their vacuum energy relative to two untwisted
lists is
\[
 h_\theta=\frac12\lim_{\varepsilon\downarrow0}
 [S_\theta(\varepsilon)+S_{1-\theta}(\varepsilon)
                      -2S_1(\varepsilon)]
 =\frac{\theta(1-\theta)}2.
\]
The divergent terms cancel before taking the limit. Zero-mode
and excitation energies are nonnegative. A scalar obeys
$h=\bar h$, so its dimension satisfies the stronger bound
\[
 \Delta\geq\max\{\theta_L(1-\theta_L),
                         \theta_R(1-\theta_R)\}.
\]
For example the last row cannot use the sum of two unequal
oscillator ground weights as a scalar dimension: the smaller
side must first rise to meet the larger. Level matching and
projection may raise the actual scalar threshold further.
Extra shift phases, or a lift of larger order, cannot lower
the oscillator floor. Consequently a scalar below $5/36$
whose two determinant signs are positive must be supported in
$K$ sectors.

\emph{Finally let the cubic select those signs.}
If $c_\phi=0$ the bound is immediate. Otherwise the complete
real doublet has the nonzero harmonic cubic already described,
with orthogonal stabilizer $D_3$. Every internal symmetry acts
within this plane by the full-eigenspace hypothesis. In particular,
the abelian quantum-character group
$\widehat{G_{\rm ab}}=\Hom(G,U(1))$ has image $1,C_2,$ or
$C_3$ in $D_3$. A character $\chi$ acts on a $g$-twisted
sector by $\chi(g)$, also with ordinary discrete torsion.
Indeed rescaling a twisted group basis by $e_g\mapsto\chi(g)e_g$
preserves its torsion-modified multiplication because
$\chi(g)\chi(h)=\chi(gh)$.

If the image is trivial, choose a unit real direction maximizing
the cubic. If it is $C_2$, choose its reflection-fixed axis.
In either case a unit Hermitian $X$ is fixed by all quantum
characters and satisfies $|C_{XXX}|=|c_\phi|/\sqrt2$.
Characters separate the abelianization, so every supporting
monodromy lies in $[G,G]$. Both determinant characters are then
positive on it. The rotation bound puts the entire support in $K$.

If the image is $C_3$, the two complex lines are simultaneous
quantum-character eigenlines. The group $D_3$ has a unique
order-three subgroup, so these are also the physical charged
lines spanned by $\phi$ and $\phi^\dagger$. Finite abelian duality identifies
their characters with inverse elements of exact order three
in $G_{\rm ab}$. Every determinant character has order at most
two and is therefore trivial on those elements. Both signs
are again positive, and the support of $\phi$ lies in $K$.
This is why low reflection-twisted sectors cannot escape the
argument by mixing with untwisted fields. We never assumed that
the physical order-three symmetry equals a quantum symmetry.

In the first two cases, lift $X$ to the intermediate Narain
theory and apply the Hermitian bound $|C_{XXX}|\leq\sqrt{2/3}$.
Its fixed-axis normalization gives
$|c_\phi|\leq\sqrt2\sqrt{2/3}=2/\sqrt3$.
In the $C_3$ case, lift $\phi$ itself and apply the complex
bound $2/\sqrt3$ directly. All lifts preserve the dimension and
normalized cubic, completing the proof.
\end{proof}

\begin{example}[A shift orbifold changes the radius]
The kernel step has a familiar one-boson model. Use the circle
convention of Chapter~\ref{ch:cft}:
\[
 h=\frac12(m/R+wR/2)^2,\qquad
 \bar h=\frac12(m/R-wR/2)^2,
 \qquad m,w\in\Z,
\]
gauge the translation $X\mapsto X+2\pi R/n$. It acts on a
momentum vertex by $e^{2\pi im/n}$, so projection requires
$m=nm'$. The $k$-twisted sector has $w\in\Z+k/n$;
combining all sectors gives $w=w'/n$ with $w'\in\Z$.
Substitution yields the same weight formulas with radius
$R'=R/n$ and integer charges $m',w'$. No oscillator has been
twisted. The smaller invariant momentum lattice and the new
fractional winding sectors restore an even self-dual charge
lattice together. Keeping only the invariant untwisted states
would omit the second half of this construction.
\end{example}

\subsection{What happens when currents are enhanced}

At a rank-two torus point, chiral charges of squared norm two form a
simply laced root system of rank at most two. The maximal
weight-one Lie algebra on either chirality is therefore one of
\[
 \mathfrak u(1)^2,\qquad
 \mathfrak{su}(2)_1\oplus\mathfrak u(1),\qquad
 \mathfrak{su}(2)_1^{\oplus2},\qquad
 \mathfrak{su}(3)_1.
\]
Their affine and Heisenberg stresses exhaust charge two. A
finite symmetry preserves the full current algebra, but may
move its chosen Cartan. This is the gap left by the preceding
proposition, not a reason to apply it without checking its
hypothesis.

\begin{constructioninput}[Completed enhanced-parent inventory]
\label{input:Narain-enhanced}
For an ordinary consistent effective finite orbifold of any
unitary rank-two Narain theory, including enhanced parents,
asymmetric chiral actions and permitted discrete torsion, an
isolated order-three scalar doublet obeys
\[
 |c_\phi|\leq2/\sqrt3
 \quad\hbox{if}\quad
 \frac{13}{100}<\Delta<\frac{61}{450}.
\]
The completed October 1 inventory combines the preceding
current-preserving theorem with specialized finite-group
classification and twisted-spectrum lemmas. Those additional
lemmas are imported here; the following case split explains
their role without claiming to reproduce their full proofs.
\end{constructioninput}

The additional inputs have a simple interpretation.
A current line preserved up to sign can be squared to give a
self-adjoint invariant operator on the isolated scalar plane.
When the sign-fixing kernel has a complementary $SU(2)$
projection with element orders at most five, its spectral
restrictions force the cubic to vanish
for $1/8<\Delta<9/50$. On an $A_2$ chirality, bounded inner
projective orders at most seven instead give no scalar at all
for $13/100<\Delta<5/36$, including actual outer extensions.
For two $A_1$ factors with bounded projections, the relevant
scalars are odd under the quantum character of factor exchange,
so their cubic vanishes in the required interval. These results
control cases where no invariant Cartan is immediately visible.

The current-algebra inventory then proceeds as follows.
\begin{itemize}
\item For $\mathfrak u(1)^2$, the whole current space already is
an invariant Heisenberg space.
\item For $\mathfrak{su}(2)\oplus\mathfrak u(1)$, the central
line is characteristic. A finite $SO(3)$ projection either has
orders at most five, invoking the current-line test, or is cyclic
or dihedral with a preserved rotation axis. That axis and the
central line give an invariant Cartan.
\item For two $\mathfrak{su}(2)$ factors, first take the kernel
of the factor permutation. Two bounded projections use the
two-factor test. Two cyclic or dihedral projections admit axes
compatible with any exchange: choose one axis and transport it
by an exchange, whose square belongs to the kernel. If just one
projection has large order, the factors cannot be exchanged;
its axis is a preserved current line and the other projection
is bounded.
\item For $\mathfrak{su}(3)$, the finite inner image can be
primitive irreducible, imprimitive, or reducible. The primitive
exceptional groups have bounded projective orders and use the
spectral test. Apart from bounded exceptions, an irreducible
imprimitive group has a unique invariant triangle of complex
lines; every normalizer preserves that triangle and hence its
Cartan. Reducible $2+1$ blocks have a characteristic central
current line. Three distinct joint-character lines give a
Cartan directly, while repeated joint characters fall under
the current-line or bounded-order cases.
\end{itemize}

Outer $SU(3)$ automorphisms are semilinear in this discussion:
they can conjugate the complex coordinates. Uniqueness of an
invariant triangle makes its preservation by such a normalizer
valid; merely choosing a convenient triangle would not suffice.
The exceptional small cases are handled by the separate bounded
spectral lemmas. Left and right current algebras may have different
enhancements. Apply the inventory to each independently, then
conjugate their invariant Cartans by the connected current
groups. This is a change to an isomorphic toroidal presentation,
after which Proposition~\ref{prop:asymmetric-doublet} applies.

The interval in Construction input~\ref{input:Narain-enhanced}
lies inside every component interval. Its upper endpoint
$61/450$ is excluded because it can occur in a bounded
two-factor spectrum. Neither the inventory nor the simpler
doublet theorem classifies all compact $c=2$ CFTs, generalized
categorical gaugings, ineffective topological factors, or
noncompact theories.

For comparison, the illustrative continuum hypothesis
$\Delta_\phi=.134$, together with a limiting cubic near the finite
estimate $|c_\phi|\simeq1.584$, would violate the bound in both
stated intervals if established.
They are not confidence intervals. Most importantly, isolation
in one measured charge sector is not isolation of the entire
continuum scalar shell. An additional exactly degenerate scalar
removes a premise even when it has a different charge. These
theorems supply conditional exclusions of construction classes,
not an identification of the chain's continuum.

\section{Ordinary crossing can demand an additional scalar}
\index{crossing symmetry!positive functional}

Crossing is useful without first guessing a named CFT. Choose the
Hermitian field $X=(\phi+\phi^\dagger)/\sqrt2$ and make the charged
self cubic positive by a field phase. Restrict its four-point
function to a line. Unitarity gives a positive diagonal
$SL(2,\R)$ block expansion with
\[
 g_D(z)=z^D\,{}_2F_1(D,D;2D;z),\qquad g_0=1.
\]
Writing $\delta=\Delta_\phi$, define crossing vectors
\[
 F_D(z)=(1-z)^{2\delta}g_D(z)-z^{2\delta}g_D(1-z).
\]
Crossing is $\sum_Da_DF_D=0$ with $a_D\geq0$. The known identity,
charged-doublet, and epsilon terms have weights
\[
 a_0=1,\qquad a_\delta=|c_\phi|^2/2,
 \qquad a_e=g_\phi^2,
\]
where $e=\Delta_\epsilon$ and
$g_\phi=C(\phi,\phi^\dagger,\epsilon)$.

A linear functional can separate allowed spectra. At $z=1/2$, put
$C=D(D-1)$, $r_D=g_D'/g_D$, and $\kappa=6/5$. The block Casimir
equation gives $g_D''=2g_D'+8Cg_D$ and
$g_D'''=8(C+1)g_D'$. Direct differentiation therefore yields
\begin{align*}
 \mathcal L(F)&=\frac{F'''-8\kappa F'}{16\,2^{-2\delta}},\\
 \frac{\mathcal L(F_D)}{g_D}
 &=(C+A-\kappa)r_D-4\delta(3C+B-\kappa),\\
 A&=1-6\delta+6\delta^2,\qquad
 B=1-3\delta+2\delta^2.
\end{align*}
Euler's integral for the hypergeometric function implies
\[
 2\sqrt2D\leq r_D\leq3D,
 \qquad (2/3)^D\leq g_D(1/2)\leq(1+2^{-D})/2.
\]
These are bounds at every positive dimension, not estimates at a
finite sample of dimensions.

To derive them, let $t$ have the symmetric beta distribution with
density proportional to $t^{D-1}(1-t)^{D-1}$ on $[0,1]$. Euler's
integral writes
\[
 g_D(1/2)=2^{-D}\mathbb E[(1-t/2)^{-D}].
\]
The weight $w(t)=(1-t/2)^{-D}$ is increasing. In the probability
measure reweighted by $w$, its mean $\mu$ therefore satisfies
$\mu\geq1/2$: the covariance of two increasing real functions is
nonnegative. Integrating the derivative of
$t^D(1-t)^D(1-t/2)^{-D}$ gives
\[
 \mathbb E_w\!\left[(2-t)-\frac2{2-t}\right]=0.
\]
Jensen's inequality for $u\mapsto1/u$ implies
$(2-\mu)^2\geq2$, hence $\mu\leq2-\sqrt2$.
Differentiating Euler's representation of $g_D$ and using the same
identity yields $r_D=4D-2D\mu$. This proves the two derivative
bounds. Finally, convexity bounds $w(\mathbb E t)$ below
$\mathbb E w$, while the chord joining $w(0)$ and $w(1)$ bounds it
above. Multiplication by $2^{-D}$ gives the two block bounds.

For $\delta\in[.13,.14]$, substituting the derivative bounds proves
$\mathcal L(F_D)>0$ for $0<D\leq.16$ and for $D\geq1.75$.
For the latter range, $C\geq21/16$ makes the coefficient of $r_D$
positive; the lower derivative bound gives a positive polynomial
lower bound. Explicitly, using $2\sqrt2>113/40$ and $D\geq7/4$,
\[
 \frac{\mathcal L(F_D)}{g_D}
 \geq\left(\frac{791}{160}-\frac{42}{25}\right)C
 -\frac{791}{160}\frac{1153}{1250}
 +\frac{36153}{125000}
 \geq .012780875>0.
\]
The coefficient of $C$ is positive, so its minimum occurs at
$C=21/16$; there is no untested high-dimension tail. For the
former range, the coefficient of $r_D$ is negative, and the
upper derivative bound gives
\[
 3D^3-4.68D^2-1.0872D+.289224,
\]
which decreases to the positive value $.007752$ at $D=.16$.

\begin{example}[A hand crossing certificate]
Assume the illustrative continuum values
\[
 \delta=.134,\quad e=.263,\quad |c_\phi|=1.584,
 \quad g_\phi=1.3.
\]
The same bounds give
$\mathcal L(F_0)>.303$,
$\mathcal L(F_\delta)>.9\times.083$,
and $\mathcal L(F_e)>-(37/40)\times.246$.
Since $|c_\phi|^2/2>5/4$, the known contribution is greater than
\[
 .303+\frac54\frac9{10}(.083)
       -\frac{169}{100}\frac{37}{40}(.246)
 =.0118155>0.
\]
All numbers in this inequality are terminating rationals. If every
other exchange lay outside $(.16,1.75)$, positivity would make the
whole functional crossing sum positive, a contradiction. Hence an
additional coupled exchange is required in that interval.
\end{example}

For identical bosonic scalars, nonzero-spin primaries have even integer
spin and hence dimension at least two by unitarity. Below $1.75$,
the required exchange is therefore scalar. One must still check
that it is not merely a new diagonal $SL(2,\R)$ primary inside an
already listed two-dimensional conformal family. For a scalar
family, the left and right global weights agree. Their tensor
product decomposes into diagonal primaries at relative levels
$n=0,1,2,\ldots$, with interchange parity $(-1)^n$. The scalar
three-point tensor is symmetric in the two equal factors, so its
odd-relative-level coefficients vanish. The first extra coupled
diagonal primary is at $\Delta+2$. A new global primary inside a
Virasoro module also first occurs at level two, since level one is
the ordinary translation descendant. Neither mechanism can produce
the required exchange below $1.75$. The conclusion is an additional
two-dimensional scalar Virasoro primary with a nonzero
$X\times X$ coupling, possibly degenerate with a supplied scalar.

The displayed numbers are hypotheses about continuum data, not
confidence intervals extracted from the chain. Moreover, the finite
chain already contains an appropriate higher charged candidate.
Its saved $L=15$ dimension in one prescribed scale is about $.731$,
with a nonzero coefficient around $.45$. Including it changes the
finite functional sum's sign. Thus the crossing demand is compatible
with an observed extra channel. Partial sums neither solve crossing
nor show a violation; descendants and the high spectrum remain.

\section{A nonreal invariant and the scope of a Cardy test}
\index{Cardy theory!phase obstruction}

The light three-point tensor can constrain more than individual OPE
magnitudes. Assume the complete real scalar shells have dimensions
$2,2,4$ for the phi, chi, and psi/eta groups, with the charge pattern
above, nonzero self cubics, and the stated physical parity. The
nonzero self cubics restrict the orthogonal transformations on the
two light planes to $D_3$. On the four-plane, the cubic is
proportional to $\Re(z^3+w^3)$; its Hessian determinant is
proportional to $|z|^2|w|^2$, forcing a symmetry to preserve or
exchange the two coordinate planes. The nonzero mixed triangles
then recover the normal $C_3\times C_3$ subgroup.

If $\Im I_\lambda\ne0$, simultaneous inversion of both charge axes
is forbidden. With the given parity, the image of all scalar-tensor
symmetries on these shells is $C_{3,\alpha}\times D_{3,S}$.
This classifies an image on selected fields, not a possible kernel
acting only on higher states, and it does not turn spatial parity
into an orientation-preserving internal line.

One consequence concerns a specific rational full-CFT construction.
Assume a unitary bosonic rational chiral algebra $V$ with positive
nonvacuum lowest weights and its canonical $A=1$ Cardy sewing,
\[
 \HH_{\rm bulk}=\bigoplus_iM_i\otimes\overline{M_i^*}.
\]
The canonical construction includes the OPE multiplication and sewing
of Chapter~\ref{ch:sewing}, not only the displayed Hilbert-space sum.
The positivity hypothesis ensures that $V$ is the maximal
holomorphic algebra of this full theory: only the right vacuum
module can contain a state of right weight zero. Every internal
symmetry preserves this maximal left algebra and its right analogue
setwise, even if it acts nontrivially on their fields. It therefore
permutes the irreducible module-pair summands. Without this step,
counting summands of a chosen smaller chiral algebra would not
constrain an arbitrary physical symmetry.

At a scalar weight $0<h<1$, all states in such a summand are
chiral ground states. A module ground space of dimension $n_i$
contributes $n_i^2$ bulk directions. An isolated real charged
doublet must therefore arise from a nonselfdual pair with
one-dimensional ground spaces: two isolated selfdual real lines
cannot carry a nontrivial order-three rotation by permutation.

The canonical identity-algebra sewing has a spatial parity
exchanging the two chiral sides. On each such doublet it conjugates
the charge. Applying this to both light doublets gives simultaneous
inversion and forces $I_\lambda$ real under the full-shell and
mixed-tensor hypotheses. A nonreal limiting value would exclude
that canonical completion.

This is not a theorem that every unitary CFT has real three-point
coefficients, nor an exclusion of every rational CFT. Other full
algebras, other sewing data, additional degenerate fields, or a
different limit lie outside its assumptions. Even a charge-conjugation
torus partition function alone is weaker than the specified canonical
sewing hypothesis.

\section{A dimension obstruction for rational lattice orbifolds}
\index{orbifold!weak integrality}

The scalar-shell arguments above need detailed low-energy information.
A different test uses only the finite category of defects, provided one
specifies the chiral algebra that those defects preserve. It extends the
pointed-lattice argument of Chapter~\ref{ch:cosets} to rational finite
orbifolds, including their twisted sectors.

Let $L$ be an even positive-definite lattice and $G$ a finite faithful
group of unitary automorphisms of $V_L$. Put $W=V_L^G$. Assume that $W$
is rational and $C_2$-cofinite, and that every irreducible $g$-twisted
$V_L$-module, except the untwisted vacuum module itself, has strictly
positive lowest weight. These are the hypotheses of the orbifold
global-dimension theorem being used; finiteness of $G$ alone is not
being offered as a proof of them. The theorem gives
\[
 \FPdim\Rep(W)=|G|^2\FPdim\Rep(V_L)=|G|^2|L^*/L|.
\]
The last equality follows because all lattice modules are invertible
and have dimension one.

\begin{proposition}[The retained rational-orbifold algebra matters]
\label{prop:rational-orbifold-Haagerup}
Assume the preceding orbifold hypotheses and a unitary rational full
CFT described over $\mathcal C=\Rep(W)$ by a simple separable Frobenius
algebra $A$, with the usual local sewing construction. Its
$W$-preserving defect category cannot contain $\mathcal H_3$ as a
full fusion subcategory.
\end{proposition}
\begin{proof}
These defects form the fusion category
$\mathcal D={}_A\mathcal C_A$. Here simplicity of $A$ means
simplicity as an $A$--$A$ bimodule:
$\operatorname{End}_{A-A}(A)=\C$, so the defect unit is simple.
The categories $\mathcal D$ and $\mathcal C$ are Morita equivalent,
so their Frobenius--Perron global dimensions agree. Therefore
$\FPdim\mathcal D=|G|^2|L^*/L|$ is an integer. A fusion category with
integer global Frobenius--Perron dimension is called
\emph{weakly integral}; the corresponding dimension theorem says
that every simple $X$ satisfies $\FPdim(X)^2\in\Z$.
The Haagerup simple $\rho$ violates this necessary condition:
\[
 \FPdim(\rho)^2=d^2=\frac{11+3\sqrt{13}}2\notin\Z.
\]
A full fusion embedding preserves simplicity and these positive
dimensions, giving a contradiction.
\end{proof}

The simplicity qualification is important: a sum of simple objects can
have irrational squared dimension even in a weakly integral category.
For a genuine unitary strong tensor action of $\mathcal H_3$, however,
Proposition~\ref{prop:c-faithful-descent} already proves full
faithfulness into a semisimple rigid target with simple unit. Thus
sending $\rho$ to a reducible defect does not evade this particular
Haagerup obstruction.

\begin{example}[A finite-group center has noninvertible objects]
Consider the untwisted group category $Z(\mathrm{Vec}_{S_3})$.
Its simples are labeled by a conjugacy class $[g]$ and an irreducible
representation $\pi$ of its centralizer; their dimensions are
$|[g]|\dim\pi$. The identity class gives dimensions $1,1,2$.
The transposition class has size three and centralizer $\Z_2$,
giving $3,3$. The three-cycle class has size two and centralizer
$\Z_3$, giving $2,2,2$. Consequently
\[
 \FPdim Z(\mathrm{Vec}_{S_3})
 =1^2+1^2+2^2+3^2+3^2+2^2+2^2+2^2=36=|S_3|^2.
\]
Finite-group constructions can therefore produce noninvertible
sectors while still obeying strong arithmetic restrictions. Weak
integrality does not require every simple dimension to be an integer:
the Ising dimensions $1,1,\sqrt2$ illustrate the distinction. Their
squares remain integers, unlike the Haagerup value above.
\end{example}

The algebra $A$ may encode a non-diagonal pairing or a local extension;
the proof is not restricted to the Cardy choice $A=\one$. It does
require the proposed defects to preserve $W$ on both chiralities.
Preserving only the total stress is weaker, and does not place a
defect in ${}_A\mathcal C_A$. Nonrational radii or orbifolds, a new
action breaking $W$, and asymmetric constructions without this
specified rational category remain outside the argument. Conversely,
within its stated framework the exclusion needs no extrapolated scalar
gap, isolated doublet, or numerical OPE coefficient.

\section{Nearby rational theories: test the retained algebra}
\index{chiral extension!candidate tests}

A central charge is a coarse invariant. The rational numbers
$21/10$, $25/12$, $13/6$, and $44/21$ all lie near the finite
Haagerup diagnostics, but a theory at one of those charges must
still have the required symmetries and scalar spectrum. We now
describe a second family of tests, complementary to the canonical
Cardy phase test: retain a specified rational chiral algebra, allow
its bosonic local extensions and physical non-diagonal pairings,
and ask whether any such full theory has enough internal symmetry.
Throughout this test assume a bosonic unitary rational full CFT with
a unique vacuum, retaining the stated algebra in both chiralities
and satisfying the usual completeness and sewing assumptions.

The distinction between a retained algebra and a number is
essential. Excluding full theories containing one particular
$W$-algebra does not exclude every theory with its central charge.
A deformation can also destroy the proposed chiral subalgebra,
even when its ultraviolet description used it.

\subsection{Why two order-three actions can be compulsory}

Suppose the ordinary internal part $S$ of translation survives
faithfully, has order three, and commutes with the full Haagerup action. It cannot
equal $\alpha$ or $\alpha^2$, because
\[
 \alpha^k\rho=\rho\alpha^{-k}
\]
and the two objects $\rho\alpha^k,\rho\alpha^{-k}$ are distinct
for $k=1,2$. Under the stated faithfulness assumptions,
$\alpha,S$ generate $\Z_3^2$. If one then gauges a suitable
order-three subgroup, a mixed anomaly can alter the extension law
and produce $\Z_9$ rather than a direct product. An obstruction
to every subgroup of order nine covers both possibilities. This
discussion assumes a precise original-versus-gauged field
dictionary; a scalar in an imposed spatially twisted sector is
not automatically a local scalar of the original theory.

\subsection{The extension and symmetry strategy}

There are several logically separate steps. Integer conformal
weight is necessary for a module to occur in a bosonic chiral
extension, because the algebra object must have trivial twist.
It is not sufficient: multiplication, associativity, commutativity,
and positivity still constrain the support. A useful additional
necessary condition is total positivity of each constituent's
quantum dimension under all number-field embeddings. Once the
actual extension algebras are classified, their chiral
automorphisms and their invertible representation objects can be
studied.

For a physical rational full CFT, pass to the maximal left and
right chiral algebras that actually occur. The full theory then
has a multiplicity-one graph pairing
\[
 \HH=\bigoplus_X X\otimes\overline{F(X)}
\]
for a braided equivalence of the corresponding representation
categories, with its associative OPE multiplication. A symmetry
fixing both chiral algebras pointwise acts by a phase $\lambda_X$
on each summand. Nonzero OPE channels require
\[
 \lambda_X\lambda_Y=\lambda_Z
 \quad\hbox{whenever }N_{XY}^{Z}>0.
\]
Thus these phases are tensor automorphisms of the identity. In a
modular category their group is identified with the invertible
objects through the nondegenerate braiding. Bounding that group
and the chiral automorphism images bounds the full internal
symmetry. The OPE assumption is important: an arbitrary formal
matrix commuting with $S,T$ need not provide this graph algebra.

\begin{example}[All minimal-model products at $c=21/10$]
The unitary Virasoro minimal charges are
$c_m=1-6/[m(m+1)]$, $m\geq3$; write $M_m=M(m,m+1)$.
A product at $21/10$ must have
exactly three nontrivial factors. Two have charge below two;
five have charge at least $5/2$. Four Ising factors give two,
and replacing any by the next model raises the charge by at
least $1/5$, overshooting $21/10$.

Order the three labels. If the smallest is at least five, the
total is at least $12/5$, which is too large. If the smallest is four, all three must
be four. If it is three, the other charges sum to $8/5$.
Taking the next label to be three, four, or five would require
the last charge to be respectively $11/10$, $9/10$, or $4/5$.
Only the last occurs, since $c_m=9/10$ would require
$m(m+1)=60$. Larger next labels already overshoot. Hence the
complete product list is
\[
 M_4^{\otimes3},\qquad M_3\otimes M_5\otimes M_5.
\]
This elementary classification is the beginning of the test,
not its conclusion; non-diagonal pairings and extensions must
still be included.
\end{example}

For three tricritical Ising factors, the only possible local
extension supports lie in the even products of their three
weight-$3/2$ fermion currents, a bosonic $\Z_2^2$. These add no
weight-two fields. Chiral automorphisms can therefore permute
the three primitive stress idempotents, with a two-group kernel.
Their order-three part is at most one three-cycle. That cycle
acts faithfully on suitable module classes, so a multiplicity-one
left-right graph pairing cannot admit independent left and right
three-cycles.

There is also a short arithmetic check on the remaining symmetry
kernel. Every symmetry in that kernel is an invertible defect
preserving the original product algebra $V=M_4^{\otimes3}$.
All $V$-preserving defects form a bimodule fusion category Morita
equivalent to $\Rep(V)$, of global dimension
$64(\varphi+2)^3$, where $\varphi=(1+\sqrt5)/2$. The smaller
symmetry kernel need not be that whole category. An invertible
order-three defect would give a fusion subcategory of dimension
three. Fusion-subcategory divisibility would force
$64(\varphi+2)^3/3$ to be an algebraic integer. Its field norm is
\[
 \frac{64^2\,5^3}{9},
\]
which is not an integer, a contradiction. Combining the chiral
and kernel arguments excludes an internal subgroup of order nine
throughout this retained-algebra family. The norm calculation is
not a statement that every category of nonintegral dimension has
no group symmetry; it uses this particular algebraic integer.

For $M_3\otimes M_5^{\otimes2}$, the only integer-weight supports
are the two weight-three $M_5$ currents and their product. Thus
the possible extensions are indexed by the five subgroups of
$\Z_2^2$. Resolving the local modules shows that order-nine
symmetry forces both $M_5$ factors to extend to Potts. The
resulting retained algebra is Ising times two Potts algebras.
It necessarily contains a unique scalar
\[
 \sigma_{\mathrm I}\otimes1\otimes1,
 \qquad \Delta=1/8,
\]
in every physical pairing. To see the last assertion directly,
Potts-pair weight residues have denominator fifteen, whereas
differences of distinct Ising weights $0,1/2,1/16$ do not.
$T$-commutation separates three Ising blocks $N_i$ of the modular
invariant. The vacuum diagonal block of $S$-commutation says
that $N_0$ commutes with the Potts-pair matrix $S_{\mathrm P}$.
The off-diagonal vacuum-row equation is
\[
 (S_{\mathrm I})_{0i}S_{\mathrm P}N_i
 =N_0S_{\mathrm P}(S_{\mathrm I})_{0i}.
\]
The nonzero Ising vacuum row and invertibility of $S_{\mathrm P}$
therefore give $N_i=N_0$: the invariant is
$I_{\mathrm I}\otimes N_0$. This retains the Ising spin field
with the vacuum Potts pair. No other scalar occurs at $1/8$, since the smallest
nonvacuum Potts chiral weight is $1/15>1/16$. A unique Hermitian
scalar is fixed by any odd-order symmetry: its real one-dimensional
orthogonal action is only a sign.

Consequently this second family conflicts with a limiting jointly
neutral scalar gap above $1/8$, if both order-three actions survive.
The chain's finite neutral gap near $.26$ motivates the test; its
limit remains an additional premise.

\subsection{Why a nearby ordinary \texorpdfstring{$W_3$}{W3} model is not enough}
\label{sec:identification-W3}

The unitary minimal algebras $W_3(p,p+1)$, $p\geq4$, have
\[
 c_p=2-\frac{24}{p(p+1)}.
\]
They approach two from below. Approximate central-charge agreement near
two therefore motivates testing this family, while an established charge
strictly above two would already distinguish it. Here we ask whether they
can support the required defects while those defects preserve this
rational $W_3$ algebra. Throughout this subsection, a full theory means
the unitary rational sewing construction from a separable Frobenius
algebra $A$ that is simple as an $A$--$A$ bimodule.

Here is a first arithmetic filter that can be understood without the
classification of full theories. Put
\[
 K_p=\Q(\zeta_{2p},\zeta_{2(p+1)}),\qquad
 \zeta_n=e^{2\pi i/n}.
\]
The positive quantum dimensions of the chiral simple objects belong to
$K_p$. In the rational sewing description by a simple separable Frobenius
algebra $A$, a defect $X$ has normalized dimension
\[
 \FPdim_{{}_A\mathcal C_A}X
   =\frac{\FPdim_{\mathcal C}X}{\FPdim_{\mathcal C}A}
   \in K_p,
 \qquad \mathcal C=\Rep\bigl(W_3(p,p+1)\bigr).
\]
The numerator and denominator use the underlying objects in
$\mathcal C$. This also applies to a reducible defect, since sums stay
in the same field.

Only primes dividing $2p(p+1)$ can ramify in this cyclotomic field.
On the other hand, $\Q(\sqrt{13})$ has discriminant thirteen, so
thirteen ramifies in it. The fact used here is that a subfield cannot
introduce a ramified prime absent from its containing field. Since
$\sqrt{13}=2d-3$, a defect of Haagerup dimension
$d=(3+\sqrt{13})/2$ would therefore require
\[
 13\mid p(p+1).
\]
The first surviving height is $p=12$, with $c_p=24/13$.
This is a necessary test, not a proof that this height works; it
still leaves infinitely many heights, with both adjacent heights
at least twelve.

\emph{An elementary gap controls ordinary affine dimensions.}
For affine $SU(3)$, let the \emph{height} $n$ mean level $n-3$.
Shifted labels are positive integers $a,b$ with $a+b<n$, and
\[
 D_n(a,b)=\frac{\sin(\pi a/n)\sin(\pi b/n)\sin(\pi(a+b)/n)}
 {\sin^2(\pi/n)\sin(2\pi/n)}.
\]
The $W_3$ dimensions are products from its two adjacent heights;
simultaneous simple-current field identification leaves them unchanged.
The vacuum label is $(1,1)$. The expression is symmetric in the
positive triple $(a,b,n-a-b)$, so sort it with
$a\leq b\leq n-a-b$. The invertible orbit has $(a,b)=(1,1)$ and
dimension one. The fundamental and its conjugate orbits have
\[
 f_n=D_n(1,2)=1+2\cos(2\pi/n),\qquad
 1+\sqrt3\leq f_n<3\quad(n\geq12).
\]
Every other sorted label has $a=1,b\geq3$, or $a,b\geq2$.
Reducing to $(1,3)$ or $(2,2)$ can only lower the dimension:
the derivative signs of the numerator with respect to $a,b$ are those
of $\sin((2a+b)\pi/n)$ and $\sin((a+2b)\pi/n)$, nonnegative in
the sorted region. Moreover, with $[m]=\sin(m\pi/n)/\sin(\pi/n)$,
\[
 \frac{D_n(2,2)}{D_n(1,3)}=\frac{[2]^2}{[3]}>1,\qquad
 D_n(1,3)=2\cos(2\pi/n)+4\cos^2(2\pi/n)\geq3+\sqrt3.
\]
Since $3<d<3+\sqrt3$, no affine simple at these heights has dimension
$d$. Two noninvertible factors have product at least
$(1+\sqrt3)^2>d$. This handles the diagonal Cardy theory, but not
yet every sewing algebra.

\begin{constructioninput}[Categorical classification inputs for this test]
The unitary diagonal affine coset has its standard full branching rule
and extension/local-module correspondence. The complete $W_3$ modular
invariant classification has exceptional affine heights $8,12,24$.
At heights at least twelve away from these exceptions, the
indecomposable module categories of the
ordinary affine $SU(3)$ fusion category are the simple-current cases
and their products with the invertible charge-conjugation module
category. Tensoring by the latter preserves the dual fusion category
up to equivalence. These are classification inputs; the reduction below
explains why they apply to the actual sewing algebras here.
\end{constructioninput}

\emph{Pass from torus matrices to actual sewing algebras.}
Choose the two adjacent heights as $n,q$, with $3\nmid q$; their
order need not be increasing. Let $\mathcal A_n$ be the affine
$SU(3)$ category at height $n$, and $\mathcal D_q$ the triality-zero
subcategory at height $q$. Triality is the central $\Z_3$ charge,
equal to $a+2b$ modulo three in shifted labels. Since $3\nmid q$,
the pointed $\Z_3$ subcategory $\mathcal P_q$ is nondegenerate and
$\mathcal A_q\simeq\mathcal P_q\boxtimes\mathcal D_q$.

The coset/extension theorem realizes the $W_3$ category as local
modules for the diagonal $\Z_3$ graph algebra in
\[
 \mathcal A_p\boxtimes\mathcal A_4
                 \boxtimes\mathcal A_{p+1}^{\mathrm{rev}}.
\]
This is an algebraic realization, not an inference from matching
$S,T$. Its reconstruction premise holds because the denominator-vacuum
branch of the numerator vacuum contains only the $W_3$ vacuum. The
branching rule reconstructs exactly the three simultaneous
simple-current translates of the vacuum. Here charge conjugation on
the reversed factor lets us write the same current exponent in all
three graph labels. Their graph multiplication
is unique up to isomorphism because $H^2(\Z_3,\C^*)=0$.

After extracting $\mathcal D_q$, take objects with trivial
$\mathcal P_q$ component. Locality uniquely determines their
$\mathcal A_4$ component from their triality. This correction changes
the braiding but not the affine associator: the pointed height-four
category has a realization with trivial associator and braiding
$c(a,b)=e^{2\pi iab/3}$ on $\Z_3$. Induction is fully
faithful on these objects, because its Hom formula sums over the three
graph translates and only the trivial translate has trivial
$\mathcal P_q$ component. The local-category dimension is
$\FPdim(\mathcal A_n)\,3\,3/3^2=\FPdim(\mathcal A_n)$, so no
additional objects remain. Consequently
\[
 \mathcal C\simeq\widetilde{\mathcal A}_n
                      \boxtimes\mathcal D_q^{\pm},
\]
where $\widetilde{\mathcal A}_n$ is monoidally equivalent to ordinary
$\mathcal A_n$, with a triality correction to its braiding; the sign
allows reversal of braiding. Neither change alters dimensions.

Let $\mathcal L$ be the Lagrangian full center of any sewing algebra.
Among the conductor-compatible heights, the classified coefficient
matrices have the form $M_n\otimes I_{\mathcal D_q}$, up to label
conjugation. Intersect $\mathcal L$ with
$\mathcal D_q\boxtimes\mathcal D_q^{\mathrm{rev}}$, interchanging these
two factors if the braiding was reversed. The intersection
theorem makes this a commutative separable algebra, also called an
\emph{\'etale algebra}, not just a subobject.
Its graph underlying object has dimension $\FPdim(\mathcal D_q)$,
so it is already Lagrangian in that factor. After condensing it,
the factor has local category $\mathrm{Vec}$: no nontrivial object
remains for a mixed gluing. The classification of connected
\'etale algebras in a product therefore factors $\mathcal L$ as an
algebra. A possible graph autoequivalence gives an invertible module
category and does not change dimensions. Thus every simple defect has
dimension $uv$, where $u$ is a simple dimension in a Morita dual of
ordinary $\mathcal A_n$, and $v$ is an ordinary dimension in
$\mathcal D_q$. This is the step that a factorized partition function
alone would miss.

\emph{All nonexceptional heights now follow from the gap.}
In the classified simple-current cases, the algebra has support
$1$ or $1+J+J^2$. A length-three orbit gives defect dimension
$D_n(a,b)$. At a fixed point the possible dimensions are
$D_n(a,b)/3$ or $D_n(a,b)$. The stabilizer of the combined left and
right action is $\Z_3$ on a free orbit, and $\Z_3^2$ at a fixed
point. Its projective irreducible representations have dimension one,
or dimension three when the latter group's commutator pairing is
nondegenerate. Dividing the underlying-object dimension by
$\FPdim A=3$ gives the stated possibilities.
A fixed point requires $3\mid n$
and $(a,b)=(n/3,n/3)$, and
\[
 \frac{D_n(n/3,n/3)}3
 =\frac{\sqrt3}{8\sin^2(\pi/n)\sin(2\pi/n)}
 \geq3+2\sqrt3>d.
\]
Every allowed $u$ or $v$ is therefore one, lies in $[1+\sqrt3,3)$,
or exceeds $d$; two noninvertible factors are too large. The affine
classification includes all module categories at these heights,
including the charge-conjugation cases, so no sewing algebra is omitted.

\emph{Only the pair $(12,13)$ requires a separate finite input.}
The other exceptional heights, eight and twenty-four, have no adjacent
height divisible by thirteen. At height twelve, all affine dimensions
lie in $\Q(\sqrt3)$: the fundamental and its dual generate the fusion
ring and both have dimension $1+\sqrt3$. Every Morita-dual simple
dimension $u$ also lies there, by the dimension-ratio formula.
At height thirteen, $v=D_{13}(a,b)\in\Q(\zeta_{13})$, whereas
\[
 \Q(\sqrt3)\cap\Q(\zeta_{13})=\Q.
\]
If $uv=d$, then $u=d/v$ belongs to both fields, so $u$ is rational
and $v$ must be a rational multiple of $d$.

\begin{constructioninput}[The exact height-thirteen comparison]
For all sixty-six pairs $a,b>0$, $a+b<13$, the three invertible
labels $(1,1),(1,11),(11,1)$ have dimension one; all other
sixty-three dimensions have degree six over $\Q$. None is therefore
a rational multiple of the quadratic number $d$.

The finite verification takes place exactly in
$\Q[z]/(1+z+\cdots+z^{12})$, with $z=\zeta_{13}$ and
$q_{13}=-z^7=e^{\pi i/13}$. It evaluates
$[m]=(q_{13}^m-q_{13}^{-m})/(q_{13}-q_{13}^{-1})$ and
$D_{13}(a,b)=[a][b][a+b]/[2]$. Galois substitution gives exact
degrees; reality and invariance under $z\mapsto z^3$ test membership
in $\Q(\sqrt{13})$. This is an exhaustive finite algebraic input,
not a comparison of approximate decimal dimensions.
The textbook imports its archived exact certificate without replaying
that calculation; Appendix~\ref{app:verification} explains this scope.
\end{constructioninput}

For example, the fundamental value is $D_{13}(1,2)=1+z+z^{-1}$.
It generates the real subfield of $\Q(\zeta_{13})$, which has degree
six; the complete finite input extends this check to every other label,
including labels outside the triality-zero subcategory.

This disposes of every sewing algebra at the exceptional pair,
including both exceptional modular-invariant types at height twelve.
The argument even excludes a reducible defect of dimension $d$.
In the nonexceptional cases every noninvertible simple has dimension
at least $1+\sqrt3>d-1$, so such an object could only be a single
noninvertible simple, already excluded, or a sum of invertibles of
integer dimension. In the exceptional case a summand with
noninvertible $v$ again exceeds $d-1$ and would have to be the whole
object; if every $v$ is invertible, the sum stays in $\Q(\sqrt3)$.

Thus a unitary strong tensor Haagerup action preserving the stated
rational $W_3$ algebra is excluded for every $p\geq4$ and every
rational sewing algebra in the stated framework. A tensor functor
preserves the positive dimension, so its image of $\rho$ would have
supplied the forbidden dimension $d$. This says nothing about
Virasoro-only defects of a $W_3$ theory or another CFT at the same
central charge. It is also different from Chapter~\ref{ch:cosets},
where an explicitly constructed irrational factor lacked an independent
spin-three vacuum generator.

\subsection{Principal \texorpdfstring{$W$}{W} algebras and an arithmetic filter}

The unitary principal type-$A$ charges are
\[
 c(N,k)=(N-1)\left[1-
 \frac{N(N+1)}{(N+k)(N+k+1)}\right],\qquad k\geq1.
\]
At $k=1$ they lie below two. For $N=2,3$ their charge is at
most two. For $N\geq7,k\geq2$,
$c(N,k)\geq c(7,2)=34/15>2.2$. Checking $N=4,5,6$
therefore leaves exactly $W_4(8,9)$ at $13/6$ and
$W_6(8,9)$ at $25/12$ in $2<c<2.2$.

The extension problem cannot be solved by looking only at the
original simple currents. In the $W_4$ coset
$(SU(4)_4\otimes SU(4)_1)/SU(4)_5$, some noncurrent
integer-weight labels initially appear possible. A typical
denominator quantum dimension among these is
\[
 d_\mu=\frac{\sin^2(4\pi/9)}
 {\sin(\pi/9)\sin(2\pi/9)}.
\]
Its Galois image multiplying sine arguments by five is negative:
the numerator remains positive, while
$\sin(5\pi/9)\sin(10\pi/9)<0$. Choosing the common-field
automorphism $113\equiv1\pmod{16}$ and
$113\equiv5\pmod{18}$ fixes the numerator $SU(4)_4$
dimensions while producing that negative denominator image.
It is an actual automorphism of the common cyclotomic field
because $\gcd(113,288)=1$.
The corresponding coset objects fail total positivity and cannot
occur in an extension algebra. This is an exact sign argument,
without a floating-point test near zero.

After that filter, the surviving extension algebras transfer to
the known $SU(4)_4$ list, including its exceptional
$SO(15)_1$ extension. All their added fields begin at weight
four, leaving one real weight-three generator. An odd-order
chiral automorphism fixes it, as it fixes the stress; these
generate the original principal algebra as a vertex algebra.
The remaining extension automorphisms and invertible modules
are two-groups. Thus every physical full theory retaining this
principal $W_4$ algebra lacks order-three internal symmetry.

\begin{constructioninput}[Completed rational-algebra tests]
The completed exact label and extension audits establish the
following restrictions, including bosonic local chiral extensions
and physical non-diagonal pairings:
\begin{center}
\begin{tabular}{p{.34\textwidth}p{.55\textwidth}}
Retained algebra & Restriction\\\hline
$M_4^{\otimes3}$, $c=21/10$
 & No internal subgroup of order nine.\\[3pt]
$M_3\otimes M_5^{\otimes2}$, $c=21/10$
 & Order-nine symmetry forces the two Potts extensions and a
 jointly neutral scalar of dimension $1/8$; $\Z_9$ is absent.\\[3pt]
$W_6(8,9)$, $c=25/12$
 & No internal subgroup of order nine.\\[3pt]
$W_4(8,9)$, $c=13/6$
 & No internal symmetry of order three.\\[3pt]
$W(E_8)(35,36)$, $c=44/21$
 & No internal symmetry of order three; every proper bosonic
 chiral extension is excluded by the exact Galois sign test.
\end{tabular}
\end{center}
The finite exhaustive label calculations and the quoted extension
classifications are inputs to this inventory. The preceding
examples explain the main methods; they do not independently
reproduce every label in those audits.
\end{constructioninput}

For $W_6$, the only nontrivial bosonic extension is by a $\Z_3$ of
integer-spin currents. Its fixed-point resolution produces no
new invertibles beyond a $\Z_2$; its possible chiral
three-cycle acts faithfully on resolved module classes. This
again prevents two independent left-right three-cycles. For the
$E_8$ principal algebra, all five nonvacuum integer-weight
supports have a negative Galois conjugate. Its distinct real
strong-generator weights $2,8,12,14,18,20,24,30$ let one
inductively fix every generator under an odd-order automorphism,
and its representation category has no nontrivial invertible
objects. These are different mechanisms leading to the stated
symmetry restrictions.

\subsection{Why the exceptional \texorpdfstring{$E_6$}{E6} branch stays separate}

The principal coset
\[
 V=\frac{(E_6)_3\otimes(E_6)_1}{(E_6)_4}
\]
also has $c=21/10$, but the previous Galois filter does not
dispose of it. Its six integer-weight supports are the vacuum,
two currents of weight eight, a conjugate pair $X,X^*$ of
weight two, and a field $Y$ of weight three. All pass the necessary
total-positivity test. Moreover a genuine exceptional affine
$(E_6)_4$ extension supplies a commutant algebra for a proposed
coset extension. Under the specified categorical transfer, the
following partial-category calculation applies. This does not
claim an independent construction of every intermediate VOA.
The branch cannot be dismissed by
assuming that all extensions are simple-current extensions.

The unextended diagonal theory and the diagonal ordinary current
extension have a scalar at $\Delta=1/10$. That is a useful test
of those pairings, not every intermediate algebra. Under the
specified exceptional commutant-algebra transfer, a partial
extension has category
\[
 \Rep(V_Y)\simeq
 \mathcal C((E_6)_3)\boxtimes\mathrm{Ising}_{5/2}.
\]
The subscript specifies the Ising modular category with fermion
weight residue $1/2$ and spin-field residue $5/16$; it is not a
claim that an ordinary critical Ising factor of charge $1/2$
appears as a tensor factor of this VOA.
Its exact twist residues imply no positive local scalar below
$7/45$ and none in $(9/40,2/5)$ in any physical pairing retaining
that extension. Under the further fully condensed branch, the
category is
\[
 (\mathrm{Fib}^{\mathrm{rev}})^{\boxtimes3}
 \boxtimes\mathrm{Ising}_{5/2},
\]
and its allowed positive scalar dimensions below one are contained
in
\[
 \{1/5,\ 9/40,\ 2/5,\ 3/5,\ 5/8\}.
\]

These constraints need only twists, once the braided category of
the retained algebra is established. A local scalar has
$h=\bar h=\Delta/2$; if $0<\Delta<1$, its actual chiral
weight lies in $(0,1/2)$ and is determined there by its residue
modulo one. Summing the Fibonacci residues
$0,3/5,1/5,4/5$ with the Ising residues $0,1/2,5/16$
gives the last list. A pairing may omit entries; it cannot add a
residue absent from the category.

As a consistency check, the partial category has topological
central charge
$78/5+5/2=181/10\equiv21/10\pmod8$.
The fully condensed category gives
$3(-14/5)+5/2=-59/10\equiv21/10\pmod8$.
Both agree with the coset charge modulo eight. This necessary
check does not turn either categorical calculation into a
uniquely identified VOA or full theory.

The charged scalar near $.13$ and neutral scalar near $.26$
would each test these exceptional branches if their limits are
established. Yet a full CFT can retain an algebra that is
mathematically extendible without containing its largest possible
extension. One cannot impose the fully condensed twist list on
every intermediate theory. Nor does equality of representation
categories identify VOAs, characters, OPEs, or chiral
automorphisms. A complete analysis of every intermediate algebra
and full pairing remains outside the completed result.

\section{A constructive ultraviolet family: three Potts order fields}

The charge pattern also suggests an explicit organizing model. Start
with three critical three-state Potts CFTs, total $c=12/5$, and their
normalized order fields $\sigma_i$ of dimension $2/15$. Preserve
\[
 \alpha:(\sigma_1,\sigma_2,\sigma_3)
       \mapsto\omega(\sigma_1,\sigma_2,\sigma_3),
\]
and
$S:(\sigma_1,\sigma_2,\sigma_3)
\mapsto(\sigma_1,\omega\sigma_2,\omega^2\sigma_3)$.
Then $\phi=\sigma_1$, $\psi=\sigma_2$, $\eta=\sigma_3$ have the
desired ordinary charges. Reflection combined with exchanging
copies two and three gives the desired parity.

The neutral scalar $W=\sigma_1\sigma_2\sigma_3$ has dimension
$2/5$. A Hermitian perturbation $uW+\bar uW^\dagger$ can have a
genuine relative phase once the three self-cubic phases are fixed.
Thermal anisotropy preserving the exchange of copies two and three
is needed to split phi from psi/eta. The relevant couplings
generated by products of $W$ must also be retained; a one-complex-
parameter action is not closed under renormalization.

The missing charge line has a simple composite basis
\[
 w_1=\sigma_2\sigma_1^\dagger,\quad
 w_2=\sigma_3\sigma_2^\dagger,\quad
 w_3=\sigma_1\sigma_3^\dagger,
\]
all of dimension $4/15$. Choose the common Potts-field phase so that
the self cubic $c_\sigma$ is real and positive. If
$\chi=\sum_ia_iw_i$ with $\sum_i|a_i|^2=1$, factorization gives
\[
 c_\chi=c_\sigma^2\sum_ia_i^3+6a_1a_2a_3,
 \qquad\lambda=a_1,\quad\kappa=a_2,\quad\nu=0.
\]
The six mixed permutations explain how a composite field can have
a large self cubic. At the ultraviolet point, however,
$|\lambda|\leq1$, and the scalar shells have extra degeneracies.
The observed finite chain tensor would require substantial
interaction effects and a demonstrated infrared splitting.

A strong-coupling clock reduction of this family has an ordinary
$C_3^2$ duality whose line dimension, if topological at a critical
point, is $\sqrt9=3$. It is not the Haagerup rho dimension, and its
fusion square is the sum of nine invertible group lines. Ordinary
self-duality therefore does not supply the missing Haagerup
junctions or prove a continuous transition. The ultraviolet family
is a concrete place to look, not an identified CFT.

\section{The remaining identification problem}

In the completed October 5 research record, the physical regular
$\mathcal H_3$ chain $H=-\sum P_\rho$ has no identified continuum CFT.
Its data suggest a long nearly critical regime, with several
central-charge diagnostics near two but substantial drift,
separated light charge shells, and a detailed nontrivial mixed
three-point tensor. They do not yet establish ultimate
masslessness, a controlled exact central charge, or continuum
OPE coefficients with error bounds.

\begin{numericalrecord}[What the latest finite comparisons add]
The $64$-state Ward runs empirically cover all four predetermined
descendant-energy windows at $L=12$ and the first three at $L=15$.
The widest $L=15$ window remains truncated. Agreement between
independent starts and recovery of known states support this
coverage, but do not certify the entire low spectrum.
Intermediate filtering improves several energy and norm
comparisons while leaving operator-dependent amplitude
discrepancies. A vector can align closely with the proposed
Ward direction and still have the wrong normalized source
amplitude; the descendant assignment must satisfy both tests.

Two invariant local directions,
\[
 V_1=\sum_jP_{\one,j},\qquad
 V_{\rm diff}=\sum_j\bigl(P_{\alpha\rho,j}-P_{\alpha^2\rho,j}\bigr),
\]
both couple strongly to the same chosen invariant-state candidate.
Their stress-subtracted diagonal responses on four fixed states
give six nonzero pairwise minors at $L=12,15$ in the stored
floating-point data, indicating two finite response directions. This tests a concrete
common-operator model at those sizes. It does not establish two
infrared primaries or two marginal couplings; subleading fields
and finite-size corrections can produce two finite response
directions even when only one field dominates in the infrared.
Separate equal-energy comparisons in the two charged families
show small decreasing catalog-wide deviations from proportionality over
$L=9,12,15$. These aggregate diagnostics do not bound each individual
pair; nonzero individual discrepancies are retained.
They support further common-field tests but prove neither an
exact operator identity nor a limiting OPE normalization.

Two stationary module-based variational trials with height-fiber
dimensions $18/54$ have visible charged height correlations
through separation $513$. One has a pronounced period-three
pattern; the other is nearly uniform at the longest measured
distances. These are direct physical insertions, so they add
information beyond auxiliary transfer eigenvalues. Both trials
are finite-capacity local optimization endpoints. The observations
do not establish an infinite-distance plateau, a ground-state
vacuum count, or a bulk gap. An ordered finite-bond approximation
can also occur when the exact chain is critical.
\end{numericalrecord}

The constructive algebraic results are valuable on their own. The
exact module tensors and their gapped parent show that symmetry
alone does not settle the phase. The $c=8$ construction has its own
finite and analytic hypotheses. The irrational chiral axis near
$2.025004294$ has a definite stress tensor but still needs the
full local completion and defect descent required for a chain
identification. Gaussian bounds, crossing functionals, and Cardy
tests can reject specified dictionaries when their premises are
established; they cannot select a unique theory by elimination
without an exhaustive classification.

To make a limiting primary-field dictionary concrete, choose a candidate full CFT and a
fixed scaling-dimension cutoff. Track the entire low-energy space below
that cutoff along a specified size sequence, including the sectors
required by the proposal. A single assignment of these spaces must
match rescaled energies, residual momenta, and symmetry actions.
Near-degenerate multiplets permit a change of basis; they do not permit
a different field identification for each observable. The vacuum and
velocity must likewise be shared.

For each proposed continuum field, choose a controlled sequence of
lattice operators and its normalization and mixing coefficients. Use
that same choice for vacuum-to-field amplitudes, matrix elements
between primary states, and descendant matrix elements. The first
group fixes the two-point normalization; the second tests OPE
coefficients; the third tests the Ward identities and descendant
assignment. Fitting an independent scale to each matrix element would
discard precisely the consistency that makes this a field dictionary.
Several microscopic recipes for the same field provide additional
checks when their separately fixed limits agree.

A correction model should then explain several observables using the
same irrelevant or running couplings and predict levels or amplitudes
not used in the fit. Example~\ref{ex:scaling-Ising-calibration} shows
that two valid finite-size calibrations can approach the same limit
from opposite sides. Example~\ref{ex:full-current-logarithms} shows
that one running coupling can produce different logarithms in different
fields. These are lessons about consistency, not correction exponents
to assign to the Haagerup chain.

A finite, mutually consistent table can support a strong numerical
identification. A proof that the chain has that scaling limit requires
further control of convergence of the state and operator data,
compatibility as the cutoff increases, and the resulting continuum
correlations. The existence of the candidate CFT is a separate
question from this convergence. Finally, the lattice lines must have
the claimed limiting action and junction spaces, with associators
equivalent to the physical $\mathcal H_3$ data. Agreement of loop dimensions
or fusion multiplicities alone does not establish that part of the
dictionary.

For a gapped interpretation with $q$ thermodynamic vacua, choose
periodic sizes compatible with the proposed spatial period and
boundary or flux sector. Along that size sequence the corresponding
target is a complete collapsing cluster,
$E_a(L)-E_0(L)\to0$ for $0\leq a<q$, separated from a positive
bulk scale
\[
 \liminf_{L\to\infty}\bigl(E_q(L)-E_0(L)\bigr)>0.
\]
The energies are ordered with multiplicity. A period-three ordered
phase, for example, can force a domain wall on rings whose length is
not divisible by three; those rings need a separate spectral analysis.
Example~\ref{ex:gapped-clock-frustration} shows this effect exactly
in a commuting clock Hamiltonian.
Connected lengths within
individual pure vacua must stay finite; a symmetric cat can retain
long-range order across the whole vacuum cluster. The exact rank-four
comparison module does not predetermine the number of vacua of the
physical projector Hamiltonian. The critical and gapped alternatives
both call for consistent state, operator, and symmetry data.

\section{Exercises}

\begin{exercise}\label{ex:identification-product}
Derive the product-theory predictions for $\Delta_\psi$,
$\lambda$, and $I_\lambda$. Explain which of them survive an
arbitrary rephasing of the three fields.

\par\smallskip\noindent{\footnotesize\normalfont
\hyperref[sol:ex:identification-product]{Solution on page~\pageref*{sol:ex:identification-product}.}}\par
\end{exercise}
\begin{exercise}\label{ex:identification-Narain}
Prove that a totally isotropic subspace of a real vector space of
signature $(2,2)$ has dimension at most two by projecting to one
positive definite factor. Explain why the strict light-dimension
inequality is needed to turn the integral pairing into zero.

\par\smallskip\noindent{\footnotesize\normalfont
\hyperref[sol:ex:identification-Narain]{Solution on page~\pageref*{sol:ex:identification-Narain}.}}\par
\end{exercise}
\begin{exercise}\label{ex:identification-triad}
Count the charge-conserving triples in the equilateral vertex
example and recover $2/\sqrt3$. Why does this not contradict the
unequal-shell theorem?

\par\smallskip\noindent{\footnotesize\normalfont
\hyperref[sol:ex:identification-triad]{Solution on page~\pageref*{sol:ex:identification-triad}.}}\par
\end{exercise}
\begin{exercise}\label{ex:identification-crossing}
Derive the functional expression from the block Casimir identities.
Check the displayed rational lower bound $.0118155$. State exactly
which unmeasured continuum assumptions are needed to apply it to
the lattice chain.

\par\smallskip\noindent{\footnotesize\normalfont
\hyperref[sol:ex:identification-crossing]{Solution on page~\pageref*{sol:ex:identification-crossing}.}}\par
\end{exercise}
\begin{exercise}\label{ex:identification-orbifold}
For asymmetric rotation fractions $(1/12,5/12)$, compute the
two oscillator ground weights. Compare their sum with the
lower bound on a scalar dimension and explain the difference.
Why is an order-two determinant character trivial on an
order-three element of the abelianization? Finally recover
the charged bound $2/\sqrt3$ from the fixed real axis and the
Hermitian parent bound, keeping the normalization factors.

\par\smallskip\noindent{\footnotesize\normalfont
\hyperref[sol:ex:identification-orbifold]{Solution on page~\pageref*{sol:ex:identification-orbifold}.}}\par
\end{exercise}
\begin{exercise}\label{ex:identification-Potts}
Derive the cubic of $\chi=\sum_ia_iw_i$ in the three-Potts example
by independent charge conservation in the three factors. Show why
a large self cubic does not imply a large mixed coefficient
$\lambda$.

\par\smallskip\noindent{\footnotesize\normalfont
\hyperref[sol:ex:identification-Potts]{Solution on page~\pageref*{sol:ex:identification-Potts}.}}\par
\end{exercise}
\begin{exercise}\label{ex:identification-norm}
Let $\varphi=(1+\sqrt5)/2$. Compute the field norms of
$\varphi+2$ and $64(\varphi+2)^3/3$. Explain the
algebraic-integer obstruction to an order-three fusion
subcategory in the tricritical Ising defect kernel. Which extra
argument is needed to control symmetries that act nontrivially
on the chiral algebra?

\par\smallskip\noindent{\footnotesize\normalfont
\hyperref[sol:ex:identification-norm]{Solution on page~\pageref*{sol:ex:identification-norm}.}}\par
\end{exercise}
\begin{exercise}\label{ex:identification-residues}
Combine the twist residues $0,3/5,1/5,4/5$ and
$0,1/2,5/16$ and recover the displayed possible scalar
dimensions strictly between zero and one. Explain why the result
is a necessary list and why it applies only when the specified
extended algebra is actually retained by the full CFT.

\par\smallskip\noindent{\footnotesize\normalfont
\hyperref[sol:ex:identification-residues]{Solution on page~\pageref*{sol:ex:identification-residues}.}}\par
\end{exercise}
